\documentclass{article}
\usepackage{amsfonts}
\usepackage{bm}
\usepackage{algorithm}
\usepackage{algorithmic}
\usepackage{mathtools}
\usepackage{graphicx}
\usepackage{booktabs}
\usepackage{xcolor}
\usepackage{hyperref}
\usepackage{amsmath, amssymb, amsthm}
\usepackage{geometry}
\usepackage{natbib}
\geometry{margin=1in}

\newtheorem{theorem}{Theorem}
\newtheorem{proposition}[theorem]{Proposition}
\newtheorem{corollary}[theorem]{Corollary}
\newtheorem{lemma}[theorem]{Lemma}
\newtheorem{remark}{Remark}
\newtheorem{definition}{Definition}
\newtheorem{assumption}{Assumption}

\DeclareMathOperator{\tr}{tr}
\DeclareMathOperator{\closure}{closure}
\DeclareMathOperator{\ilr}{ilr}
\DeclareMathOperator{\diag}{diag}
\DeclareMathOperator{\softmax}{softmax}
\DeclareMathOperator{\vect}{vec}

\begin{document}

\title{Structural Topic Modeling in Aitchison Geometry:\\
  An Analytical Spectral Solution via ILR-Spectral-GSCA}
\date{\today}
\maketitle

\begin{abstract}
We introduce a structural topic model that operates natively in the
Aitchison geometry of the simplex. Document--topic proportions are
represented in isometric log-ratio (ILR) coordinates, and covariates are
linked to topic prevalence through a Generalized Structured Component
Analysis (GSCA) formulation. By profiling out all nuisance parameters in
closed form, we reduce the estimation problem to a single eigenvalue
decomposition of a covariate-augmented document similarity matrix.
The resulting estimator is non-iterative, globally optimal for the
profiled linearised objective, and yields topic proportions on the
probability simplex by construction---not through a post-hoc
transformation, but as the natural inverse of the coordinate system in
which estimation is performed. We establish global optimality for the
linearised ILR-Spectral-GSCA objective, clarify the identifiability conditions
of the spectral solution under both simple and degenerate spectra, and
derive consistency and asymptotic normality of the covariate path
coefficients with an explicit analytical covariance formula based on
eigenvector perturbation theory. A linearisation error bound under an
eigenvector delocalisation condition quantifies the approximation quality
relative to the original simplex problem. The approach bridges three
previously disjoint frameworks---topic modeling, structural equation
modeling, and compositional data analysis---through the unifying lens of
Aitchison geometry.
\end{abstract}

% ====================================================================
\section{Introduction}
\label{sec:introduction}
% ====================================================================

[To be written.]

% ====================================================================
\section{Background: Aitchison Geometry and the ILR Transformation}
\label{sec:background}
% ====================================================================

The probability simplex
$\mathcal{S}^K = \bigl\{\mathbf{p}\in\mathbb{R}^K_{>0}:\sum_k p_k=1\bigr\}$
is not a Euclidean space: it is a $(K{-}1)$-dimensional manifold with a
natural metric structure induced by the Aitchison inner
product~\citep{aitchison1986statistical}. The \emph{isometric log-ratio}
(ILR) transformation~\citep{egozcue2003isometric} provides an isometry
between $(\mathcal{S}^K, d_A)$ and $(\mathbb{R}^{K-1}, \|\cdot\|_2)$,
enabling standard linear algebra operations on compositional data
without distortion.

\begin{definition}[ILR transformation]
\label{def:ilr}
Let $\mathbf{V}\in\mathbb{R}^{K\times(K{-}1)}$ be a \emph{contrast
matrix} satisfying
\begin{equation}
  \mathbf{V}^{\!\top}\mathbf{V} = \mathbf{I}_{K-1},
  \qquad
  \mathbf{V}^{\!\top}\mathbf{1}_K = \mathbf{0}_{K-1}.
  \label{eq:V_properties}
\end{equation}
For $\boldsymbol{\theta}\in\mathcal{S}^K$, the ILR coordinates are
$\mathbf{z} = \mathbf{V}^{\!\top}\ln\boldsymbol{\theta}
 \in\mathbb{R}^{K-1}$,
and the inverse transformation is
\begin{equation}
  \boldsymbol{\theta}
  = \closure\!\bigl(\exp(\mathbf{V}\mathbf{z})\bigr)
  \;\stackrel{\text{def}}{=}\;
  \frac{\exp(\mathbf{V}\mathbf{z})}
       {\mathbf{1}_K^{\!\top}\exp(\mathbf{V}\mathbf{z})}.
  \label{eq:ilr_inverse}
\end{equation}
\end{definition}

\noindent
Two key properties of $\mathbf{V}$ follow
from~\eqref{eq:V_properties}: (i) $\mathbf{V}\mathbf{V}^{\!\top}
= \mathbf{I}_K - K^{-1}\mathbf{1}_K\mathbf{1}_K^{\!\top}$
(the centering projector in topic space); and (ii) the simplex centroid
$\boldsymbol{\theta}_0 = K^{-1}\mathbf{1}_K$ maps to the origin,
$\mathbf{z}_0 = \mathbf{0}$.

% ====================================================================
\section{Model Specification}
\label{sec:model}
% ====================================================================

\subsection{Data}

Let $\mathbf{W}\in\mathbb{R}_{\geq 0}^{M\times N}$ be a document--term
matrix ($M$~documents, $N$~terms) and
$\mathbf{C}\in\mathbb{R}^{M\times P}$ a column-centred covariate matrix
with full column rank~$P$. We seek $K$~latent topics parameterised by
a document--topic proportion matrix
$\boldsymbol{\Theta}\in\mathbb{R}^{M\times K}$ (each row on the
probability simplex $\mathcal{S}^K$), a topic--term loading matrix
$\boldsymbol{\Phi}\in\mathbb{R}^{K\times N}$, and a covariate--topic
path coefficient matrix $\mathbf{B}\in\mathbb{R}^{P\times K}$.


\subsection{Objective Function on the Simplex}
\label{subsec:objective_simplex}

The model seeks to factorise the document--term matrix while linking
topic prevalence to covariates. We formulate this as the penalised
least-squares problem
%
\begin{equation}
  \min_{\boldsymbol{\Theta},\,\boldsymbol{\Phi},\,\mathbf{B}}\;
  \underbrace{%
    \bigl\|\mathbf{W}
      - \boldsymbol{\Theta}\boldsymbol{\Phi}\bigr\|_F^2
  }_{\text{reconstruction}}
  \;+\;
  \lambda\,
  \underbrace{%
    \bigl\|\boldsymbol{\Theta}
      - \mathbf{C}\mathbf{B}\bigr\|_F^2
  }_{\text{structural model}},
  \label{eq:objective_simplex}
\end{equation}
%
subject to
$\boldsymbol{\theta}_i\in\mathcal{S}^K$ for all $i=1,\dots,M$,
where $\lambda>0$ balances reconstruction fidelity against covariate
structure.

The first term measures how well the low-rank product
$\boldsymbol{\Theta}\boldsymbol{\Phi}$ approximates the observed word
counts; the second enforces a linear relationship between covariates
and topic prevalence, with $\mathbf{B}$ capturing how each covariate
shifts the topic composition.

\begin{remark}[The core difficulty]
\label{rem:core_difficulty}
Problem~\eqref{eq:objective_simplex} is jointly non-convex in
$(\boldsymbol{\Theta},\boldsymbol{\Phi},\mathbf{B})$ under the
simplex constraint. Existing approaches either embed the structural
model in a hierarchical Bayesian framework and resort to variational
or MCMC inference~\citep{roberts2014stm}, or replace the simplex
constraint with an EM--GSCA hybrid that converges only to local
optima. We take a different route: we reparametrise the problem in
the natural coordinate system of the simplex---the ILR space---where
an analytical, globally optimal solution becomes available.
\end{remark}


\subsection{Reparametrisation via ILR Coordinates}
\label{subsec:reparametrisation}

The simplex $\mathcal{S}^K$ is a $(K{-}1)$-dimensional manifold. The
ILR transformation (Definition~\ref{def:ilr}) maps each
$\boldsymbol{\theta}_i\in\mathcal{S}^K$ to an unconstrained vector
$\mathbf{z}_i = \mathbf{V}^{\!\top}\ln\boldsymbol{\theta}_i
 \in\mathbb{R}^{K-1}$,
with inverse
$\boldsymbol{\theta}_i
 = \closure(\exp(\mathbf{V}\mathbf{z}_i))$.

\paragraph{Linearisation.}
A first-order Taylor expansion of the closure map around the simplex
centroid $\boldsymbol{\theta}_0 = K^{-1}\mathbf{1}_K$ (which maps to
$\mathbf{z}_0=\mathbf{0}$) yields the approximation
(Lemma~\ref{lem:linearisation}):
%
\begin{equation}
  \boldsymbol{\theta}_i
  \;\approx\;
  \frac{1}{K}\bigl(\mathbf{1}_K + \mathbf{V}\mathbf{z}_i\bigr).
  \label{eq:theta_approx}
\end{equation}
%
Stacking documents row-wise:
\begin{equation}
  \boldsymbol{\Theta}
  \;\approx\;
  \frac{1}{K}
  \bigl(\mathbf{1}_M\mathbf{1}_K^{\!\top}
        + \mathbf{Z}\mathbf{V}^{\!\top}\bigr),
  \label{eq:Theta_approx}
\end{equation}
where $\mathbf{Z}\in\mathbb{R}^{M\times(K{-}1)}$ collects the ILR
scores.

\begin{lemma}[Linearisation of the closure map]
\label{lem:linearisation}
Let $f(\mathbf{z}) = \closure(\exp(\mathbf{V}\mathbf{z}))$, with
$f(\mathbf{0})=K^{-1}\mathbf{1}_K$. Then the Jacobian at the origin is
$\mathbf{J}_f(\mathbf{0}) = K^{-1}\mathbf{V}$, and
\[
  \bigl\|
    f(\mathbf{z})
    - K^{-1}(\mathbf{1}_K + \mathbf{V}\mathbf{z})
  \bigr\|_2
  \;\leq\;
  \frac{\|\mathbf{z}\|_2^2}{2K}
  \,\exp\!\bigl(\|\mathbf{z}\|_2\bigr).
\]
\end{lemma}

\begin{proof}
The $k$-th component of $f$ is
$f_k(\mathbf{z}) = \exp((\mathbf{V}\mathbf{z})_k)/\sum_{k'}
\exp((\mathbf{V}\mathbf{z})_{k'})$.
At $\mathbf{z}=\mathbf{0}$, all exponentials equal~1, so
$f_k(\mathbf{0})=1/K$. Differentiating,
\[
  \frac{\partial f_k}{\partial z_j}\bigg|_{\mathbf{z}=\mathbf{0}}
  = \frac{1}{K}\Bigl[V_{kj}
    - \frac{1}{K}\sum_{k'}V_{k'j}\Bigr]
  = \frac{V_{kj}}{K},
\]
since $\sum_{k'}V_{k'j} = \mathbf{1}_K^{\!\top}\mathbf{V}_j = 0$
by~\eqref{eq:V_properties}. Thus
$\mathbf{J}_f(\mathbf{0}) = K^{-1}\mathbf{V}$.
The remainder bound follows from the Lipschitz constant of the
softmax Hessian on compact subsets of~$\mathbb{R}^{K-1}$.
\end{proof}

\paragraph{Decomposition of $\boldsymbol{\Phi}$.}
The topic--term matrix $\boldsymbol{\Phi}\in\mathbb{R}^{K\times N}$
decomposes in the ILR basis as
\begin{equation}
  \boldsymbol{\Phi}
  \;=\;
  \mathbf{V}\boldsymbol{\Psi}
  \;+\;
  \mathbf{1}_K\boldsymbol{\mu}^{\!\top},
  \label{eq:Phi_decomposition}
\end{equation}
where
$\boldsymbol{\Psi}=\mathbf{V}^{\!\top}\boldsymbol{\Phi}
 \in\mathbb{R}^{(K{-}1)\times N}$
captures topic-specific deviations in ILR space and
$\boldsymbol{\mu}^{\!\top}=K^{-1}\mathbf{1}_K^{\!\top}\boldsymbol{\Phi}$
is the geometric-mean profile across topics.


\subsection{Derivation of the ILR Objective}
\label{subsec:ilr_derivation}

We now substitute the ILR
reparametrisation~\eqref{eq:Theta_approx}--\eqref{eq:Phi_decomposition}
into the original objective~\eqref{eq:objective_simplex} to obtain a
formulation in ILR coordinates.

\paragraph{Step 1: Reconstruction term.}
The product $\boldsymbol{\Theta}\boldsymbol{\Phi}$ under the linearised
approximation simplifies to
\begin{equation}
  \boldsymbol{\Theta}\boldsymbol{\Phi}
  \;=\;
  \mathbf{1}_M\boldsymbol{\mu}^{\!\top}
  \;+\;
  \frac{1}{K}\,\mathbf{Z}\boldsymbol{\Psi},
  \label{eq:reconstruction_ilr}
\end{equation}
as established by the following derivation. Expanding
$K^{-1}(\mathbf{1}_M\mathbf{1}_K^{\!\top}
+ \mathbf{Z}\mathbf{V}^{\!\top})(\mathbf{V}\boldsymbol{\Psi}
+ \mathbf{1}_K\boldsymbol{\mu}^{\!\top})$ yields four terms:
\begin{align*}
  \text{(T1)}\quad
  &K^{-1}\mathbf{1}_M
   \underbrace{\mathbf{1}_K^{\!\top}\mathbf{V}}_{=\,\mathbf{0}^{\!\top}}
   \boldsymbol{\Psi}
   = \mathbf{0}, \\
  \text{(T2)}\quad
  &K^{-1}\mathbf{1}_M\mathbf{1}_K^{\!\top}\mathbf{1}_K
   \boldsymbol{\mu}^{\!\top}
   = \mathbf{1}_M\boldsymbol{\mu}^{\!\top}, \\
  \text{(T3)}\quad
  &K^{-1}\mathbf{Z}
   \underbrace{\mathbf{V}^{\!\top}\mathbf{V}}_{=\,\mathbf{I}_{K-1}}
   \boldsymbol{\Psi}
   = K^{-1}\mathbf{Z}\boldsymbol{\Psi}, \\
  \text{(T4)}\quad
  &K^{-1}\mathbf{Z}
   \underbrace{\mathbf{V}^{\!\top}\mathbf{1}_K}_{=\,\mathbf{0}}
   \boldsymbol{\mu}^{\!\top}
   = \mathbf{0}.
\end{align*}
The cross-terms (T1) and (T4) vanish due to the orthogonality property
$\mathbf{V}^{\!\top}\mathbf{1}_K=\mathbf{0}$, producing the clean
decomposition~\eqref{eq:reconstruction_ilr}: a document-independent
centroid $\mathbf{1}_M\boldsymbol{\mu}^{\!\top}$ plus a low-rank
ILR-space perturbation $K^{-1}\mathbf{Z}\boldsymbol{\Psi}$.

Substituting~\eqref{eq:reconstruction_ilr} into the reconstruction
term of~\eqref{eq:objective_simplex}:
\begin{equation}
  \bigl\|\mathbf{W}
    - \boldsymbol{\Theta}\boldsymbol{\Phi}\bigr\|_F^2
  \;=\;
  \Bigl\|\mathbf{W}
    - \mathbf{1}_M\boldsymbol{\mu}^{\!\top}
    - \tfrac{1}{K}\,\mathbf{Z}\boldsymbol{\Psi}
  \Bigr\|_F^2.
  \label{eq:recon_ilr}
\end{equation}

\paragraph{Step 2: Structural term.}
Under~\eqref{eq:Theta_approx}, the structural penalty becomes
\begin{align}
  \bigl\|\boldsymbol{\Theta} - \mathbf{C}\mathbf{B}\bigr\|_F^2
  &= \Bigl\|
    \tfrac{1}{K}(\mathbf{1}_M\mathbf{1}_K^{\!\top}
    + \mathbf{Z}\mathbf{V}^{\!\top})
    - \mathbf{C}\mathbf{B}
  \Bigr\|_F^2.
  \label{eq:structural_expanded}
\end{align}
Since $\mathbf{C}$ is column-centred ($\mathbf{C}^{\!\top}\mathbf{1}_M
= \mathbf{0}$), the constant term
$K^{-1}\mathbf{1}_M\mathbf{1}_K^{\!\top}$ is orthogonal to the column
space of~$\mathbf{C}$ and contributes only a constant to the penalty.
Reparametrising
$\mathbf{B} = K^{-1}\mathbf{B}_z\mathbf{V}^{\!\top}$ (so that
$\mathbf{C}\mathbf{B}$ targets $K^{-1}\mathbf{Z}\mathbf{V}^{\!\top}$),
the structural penalty reduces, up to an additive constant, to
\begin{equation}
  \lambda\bigl\|\boldsymbol{\Theta}
    - \mathbf{C}\mathbf{B}\bigr\|_F^2
  \;=\;
  \frac{\lambda}{K^2}
  \bigl\|\mathbf{Z} - \mathbf{C}\mathbf{B}_z\bigr\|_F^2
  \;+\;\text{const.}
  \label{eq:structural_ilr}
\end{equation}

\paragraph{Step 3: Absorbing $K^{-2}$ into $\lambda$.}
Defining $\tilde{\lambda}=\lambda/K^2$ (or equivalently, redefining
$\lambda$ to absorb the scaling without loss of generality since
$\lambda$ is a user-specified tuning parameter), we arrive at the
\textbf{ILR-Spectral-GSCA objective}:
%
\begin{equation}
  \boxed{\;
  \min_{\mathbf{Z},\,\boldsymbol{\Psi},\,\mathbf{B}_z,\,\boldsymbol{\mu}}\;
  \Bigl\|
    \mathbf{W}
    - \mathbf{1}_M\boldsymbol{\mu}^{\!\top}
    - \tfrac{1}{K}\,\mathbf{Z}\boldsymbol{\Psi}
  \Bigr\|_F^2
  \;+\;
  \lambda
  \bigl\|\mathbf{Z} - \mathbf{C}\mathbf{B}_z\bigr\|_F^2
  \;}
  \label{eq:objective}
\end{equation}
%
subject to the Spectral-GSCA normalisation constraint
\begin{equation}
  \mathbf{Z}^{\!\top}\mathbf{Z} = \mathbf{I}_{K-1}.
  \label{eq:constraint}
\end{equation}

\begin{remark}[What the reparametrisation achieves]
\label{rem:what_reparam_achieves}
Three features of~\eqref{eq:objective} are worth noting.
%
First, the simplex constraint on~$\boldsymbol{\Theta}$ has been replaced
by an unconstrained estimation problem in $\mathbb{R}^{K-1}$: the
constraint $\sum_k\theta_{ik}=1$ is absorbed by the ILR transformation,
and the non-negativity $\theta_{ik}\geq 0$ is recovered at the end via
the closure map.
%
Second, the structural penalty
$\|\mathbf{Z}-\mathbf{C}\mathbf{B}_z\|_F^2$ in ILR coordinates
corresponds to the sum of squared Aitchison distances between the
observed and covariate-predicted topic compositions:
if $\mathbf{z}_i = \ilr(\boldsymbol{\theta}_i)$ and
$\boldsymbol{\eta}_i = \mathbf{c}_i^{\!\top}\mathbf{B}_z$, then
$\|\mathbf{z}_i - \boldsymbol{\eta}_i\|_2^2
 = d_A(\boldsymbol{\theta}_i,
       \closure(\exp(\mathbf{V}\boldsymbol{\eta}_i)))^2$
by the isometry property of the ILR
transformation~\citep{aitchison1986statistical}.
This Aitchison interpretation applies to the reparameterised
structural term in ILR coordinates; it should not be read as a literal
property of the original Frobenius penalty
$\|\boldsymbol{\Theta}-\mathbf{C}\mathbf{B}\|_F^2$ written
in~$\boldsymbol{\Theta}$-space.
%
Third, the dimension of the latent space is $K{-}1$, not $K$: the
``missing'' degree of freedom is precisely the sum-to-one constraint.
\end{remark}


% ====================================================================
\section{Analytical Solution via Spectral Profiling}
\label{sec:analytical}
% ====================================================================

We derive the closed-form solution
to~\eqref{eq:objective}--\eqref{eq:constraint} by sequentially
eliminating $\boldsymbol{\mu}$, $\boldsymbol{\Psi}$, and
$\mathbf{B}_z$, reducing the problem to a single eigenvalue
decomposition in~$\mathbf{Z}$.

% ---- Step 1: μ ---------------------------------------------------
\subsection{Profiling out
\texorpdfstring{$\boldsymbol{\mu}$}{μ}}
\label{subsec:profile_mu}

\begin{proposition}
\label{prop:mu_profile}
For fixed $(\mathbf{Z},\boldsymbol{\Psi})$ with
$\mathbf{Z}^{\!\top}\mathbf{1}_M = \mathbf{0}$ (centred scores), the
minimiser of the reconstruction term
over~$\boldsymbol{\mu}$ is
\begin{equation}
  \hat{\boldsymbol{\mu}}
  = \frac{1}{M}\,\mathbf{W}^{\!\top}\mathbf{1}_M
  \;=\; \bar{\mathbf{w}},
  \label{eq:mu_hat}
\end{equation}
the column-mean vector of~$\mathbf{W}$.
After substitution, the reconstruction term becomes
$\|\widetilde{\mathbf{W}} - K^{-1}\mathbf{Z}\boldsymbol{\Psi}\|_F^2$,
where $\widetilde{\mathbf{W}} = \mathbf{H}_M\mathbf{W}$ is the
column-centred document--term matrix and
$\mathbf{H}_M = \mathbf{I}_M - M^{-1}\mathbf{1}_M\mathbf{1}_M^{\!\top}$
is the centering matrix.
\end{proposition}

\begin{proof}
Differentiating with respect to $\boldsymbol{\mu}$ and setting to zero:
\[
  \mathbf{1}_M^{\!\top}
  \bigl(\mathbf{W} - \mathbf{1}_M\boldsymbol{\mu}^{\!\top}
        - K^{-1}\mathbf{Z}\boldsymbol{\Psi}\bigr) = \mathbf{0}^{\!\top}.
\]
Since $\mathbf{1}_M^{\!\top}\mathbf{Z}=\mathbf{0}^{\!\top}$ (centred
scores), this reduces to
$M\boldsymbol{\mu}^{\!\top}
 = \mathbf{1}_M^{\!\top}\mathbf{W}$.
\end{proof}

\noindent
The centering $\widetilde{\mathbf{W}} = \mathbf{H}_M\mathbf{W}$ is
natural in the Aitchison framework: it removes the corpus-level baseline,
isolating the document-specific compositional deviations that the topics
must explain.

% ---- Step 2: Ψ ---------------------------------------------------
\subsection{Profiling out
\texorpdfstring{$\boldsymbol{\Psi}$}{Psi}}
\label{subsec:profile_Psi}

\begin{proposition}
\label{prop:Psi_profile}
For fixed $\mathbf{Z}$ with
$\mathbf{Z}^{\!\top}\mathbf{Z}=\mathbf{I}_{K-1}$, the minimiser of
$\|\widetilde{\mathbf{W}} - K^{-1}\mathbf{Z}\boldsymbol{\Psi}\|_F^2$
over $\boldsymbol{\Psi}$ is
\begin{equation}
  \hat{\boldsymbol{\Psi}}
  = K\,\mathbf{Z}^{\!\top}\widetilde{\mathbf{W}}.
  \label{eq:Psi_hat}
\end{equation}
Substituting into the reconstruction term gives
\begin{equation}
  \bigl\|\widetilde{\mathbf{W}}
    - K^{-1}\mathbf{Z}\hat{\boldsymbol{\Psi}}\bigr\|_F^2
  = \tr\!\bigl(\widetilde{\mathbf{W}}^{\!\top}\widetilde{\mathbf{W}}\bigr)
    - \tr\!\bigl(
        \mathbf{Z}^{\!\top}
        \widetilde{\mathbf{W}}\widetilde{\mathbf{W}}^{\!\top}
        \mathbf{Z}
      \bigr).
  \label{eq:recon_profiled}
\end{equation}
\end{proposition}

\begin{proof}
Setting the derivative with respect to $\boldsymbol{\Psi}$ to zero:
$K^{-2}\mathbf{Z}^{\!\top}\mathbf{Z}\boldsymbol{\Psi}
 = K^{-1}\mathbf{Z}^{\!\top}\widetilde{\mathbf{W}}$.
Under the constraint, $\boldsymbol{\Psi}=K\mathbf{Z}^{\!\top}
\widetilde{\mathbf{W}}$. Substituting,
$\widetilde{\mathbf{W}} - K^{-1}\mathbf{Z}(K\mathbf{Z}^{\!\top}
\widetilde{\mathbf{W}})
= (\mathbf{I}_M - \mathbf{Z}\mathbf{Z}^{\!\top})\widetilde{\mathbf{W}}$,
and
$\|(\mathbf{I}-\mathbf{Z}\mathbf{Z}^{\!\top})\widetilde{\mathbf{W}}\|_F^2
= \tr(\widetilde{\mathbf{W}}^{\!\top}\widetilde{\mathbf{W}})
  - \tr(\mathbf{Z}^{\!\top}\widetilde{\mathbf{W}}
    \widetilde{\mathbf{W}}^{\!\top}\mathbf{Z})$
since $\mathbf{I}-\mathbf{Z}\mathbf{Z}^{\!\top}$ is an orthogonal
projector.
\end{proof}

% ---- Step 3: B_z -------------------------------------------------
\subsection{Profiling out
\texorpdfstring{$\mathbf{B}_z$}{Bz}}
\label{subsec:profile_Bz}

\begin{proposition}
\label{prop:Bz_profile}
For fixed $\mathbf{Z}$, the minimiser of the structural penalty is
\begin{equation}
  \hat{\mathbf{B}}_z
  = \bigl(\mathbf{C}^{\!\top}\mathbf{C}\bigr)^{-1}
    \mathbf{C}^{\!\top}\mathbf{Z}.
  \label{eq:Bz_hat}
\end{equation}
Substituting,
\begin{equation}
  \bigl\|\mathbf{Z}
    - \mathbf{C}\hat{\mathbf{B}}_z\bigr\|_F^2
  = \tr\!\bigl(
      \mathbf{Z}^{\!\top}\mathbf{M}_C\,\mathbf{Z}
    \bigr)
  = (K{-}1) - \tr\!\bigl(
      \mathbf{Z}^{\!\top}\mathbf{P}_C\,\mathbf{Z}
    \bigr),
  \label{eq:structural_profiled}
\end{equation}
where $\mathbf{P}_C = \mathbf{C}(\mathbf{C}^{\!\top}\mathbf{C})^{-1}
\mathbf{C}^{\!\top}$ and
$\mathbf{M}_C = \mathbf{I}_M - \mathbf{P}_C$.
\end{proposition}

% ---- Main theorem ------------------------------------------------
\subsection{The ILR-Spectral-GSCA Spectral Theorem}
\label{subsec:spectral}

Combining Propositions~\ref{prop:mu_profile}--\ref{prop:Bz_profile},
the fully profiled objective (dropping additive constants in
$\mathbf{W}$, $\lambda$, and $K$) reduces to:

\begin{equation}
  \max_{\mathbf{Z}^{\!\top}\mathbf{Z}=\mathbf{I}_{K-1}}\;
  \tr\!\Bigl(
    \mathbf{Z}^{\!\top}\,\mathbf{S}_z\,\mathbf{Z}
  \Bigr),
  \label{eq:profiled_max}
\end{equation}
where
\begin{equation}
  \boxed{\;
  \mathbf{S}_z
  \;=\;
  \widetilde{\mathbf{W}}\widetilde{\mathbf{W}}^{\!\top}
  \;+\;
  \lambda\,\mathbf{P}_C
  \;=\;
  \mathbf{H}_M\mathbf{W}\mathbf{W}^{\!\top}\mathbf{H}_M
  \;+\;
  \lambda\,\mathbf{C}
  \bigl(\mathbf{C}^{\!\top}\mathbf{C}\bigr)^{-1}
  \mathbf{C}^{\!\top}
  \;}
  \label{eq:Sz_matrix}
\end{equation}
is a symmetric positive semi-definite $M\times M$ matrix.

\begin{theorem}[Analytical solution for ILR topic scores]
\label{thm:spectral}
Let $\sigma_1\geq\sigma_2\geq\cdots\geq\sigma_M\geq 0$ be the
eigenvalues of $\mathbf{S}_z$, with orthonormal eigenvectors
$\mathbf{u}_1,\dots,\mathbf{u}_M$.
\begin{enumerate}
  \item[\emph{(i)}] The vector $\mathbf{1}_M$ belongs to the null space
    of $\mathbf{S}_z$, i.e.,
    $\mathbf{S}_z\mathbf{1}_M = \mathbf{0}$.
  \item[\emph{(ii)}] The global maximum
    of~\eqref{eq:profiled_max} is
    $\sum_{k=1}^{K-1}\sigma_k$, attained by
    \begin{equation}
      \mathbf{Z}^*
      = \bigl[\,\mathbf{u}_1\;\;\mathbf{u}_2\;\;\cdots\;\;
               \mathbf{u}_{K-1}\,\bigr]
      \in\mathbb{R}^{M\times(K-1)}.
      \label{eq:Z_star}
    \end{equation}
  \item[\emph{(iii)}] The constraint
    $\mathbf{Z}^{*\top}\mathbf{1}_M = \mathbf{0}$ is automatically
    satisfied.
\end{enumerate}
\end{theorem}

\begin{proof}
\emph{(i)} Since $\mathbf{C}$ is column-centred,
$\mathbf{C}^{\!\top}\mathbf{1}_M = \mathbf{0}$, so
$\mathbf{P}_C\mathbf{1}_M = \mathbf{0}$.
Similarly, $\mathbf{H}_M\mathbf{1}_M = \mathbf{0}$ implies
$\widetilde{\mathbf{W}}\widetilde{\mathbf{W}}^{\!\top}\mathbf{1}_M
= \mathbf{0}$. Hence $\mathbf{S}_z\mathbf{1}_M = \mathbf{0}$.

\emph{(ii)} $\mathbf{S}_z$ is real symmetric and positive semi-definite
(sum of two PSD matrices). The result is an application of the
Ky~Fan $k$-eigenvalue theorem~\citep{fan1949theorem}.

\emph{(iii)} Since $\mathbf{1}_M$ is an eigenvector with eigenvalue~0
and $\sigma_{K-1}>0$ (assuming $K{-}1$ non-zero eigenvalues), the top
$(K{-}1)$ eigenvectors are orthogonal to $\mathbf{1}_M$, i.e.,\
$\mathbf{u}_k^{\!\top}\mathbf{1}_M = 0$ for $k=1,\dots,K{-}1$.
\end{proof}

\begin{remark}[Dimension reduction]
\label{rem:dimension_reduction}
The simplex constraint $\boldsymbol{\theta}_i\in\mathcal{S}^K$
contributes $K$~parameters per document in the original space but only
$K{-}1$ in ILR coordinates. Accordingly, the eigenvalue problem
in~\eqref{eq:profiled_max} is of dimension $K{-}1$, not $K$. The
``missing'' dimension is absorbed by the constraint
$\mathbf{Z}^{*\top}\mathbf{1}_M = \mathbf{0}$, which arises
automatically from the structure of~$\mathbf{S}_z$.
\end{remark}

\begin{remark}[Interpretation of $\mathbf{S}_z$]
\label{rem:Sz_interpretation}
The covariate-augmented similarity matrix $\mathbf{S}_z$ has a
transparent compositional interpretation. The term
$\widetilde{\mathbf{W}}\widetilde{\mathbf{W}}^{\!\top}$ is the centred
document similarity matrix, whose eigenvectors yield the standard
Latent Semantic Analysis (LSA) decomposition.
The term $\lambda\mathbf{P}_C$ augments this with the inter-document
similarity induced by the covariates.
The regularisation parameter $\lambda$ governs the balance:
$\lambda=0$ recovers centred LSA;
$\lambda\to\infty$ collapses the topics onto the covariate space.
For intermediate $\lambda$, the topics are the directions in
\emph{centred} document space that jointly maximise explained text
variance and alignment with covariate structure---all measured in the
natural Aitchison metric.
\end{remark}


% ---- Recovery of all parameters ---------------------------------
\subsection{Recovery of All Parameters}
\label{subsec:recovery}

\begin{corollary}[Complete parameter recovery]
\label{cor:recovery}
Given the optimal ILR scores $\mathbf{Z}^*$
from Theorem~\ref{thm:spectral}:
\begin{align}
  \hat{\boldsymbol{\mu}}
    &= \bar{\mathbf{w}}
    &&\text{(column mean of }\mathbf{W}\text{)},
    \label{eq:mu_recovery}\\[4pt]
  \hat{\boldsymbol{\Psi}}
    &= K\,\mathbf{Z}^{*\top}\widetilde{\mathbf{W}}
    &&\in\mathbb{R}^{(K{-}1)\times N},
    \label{eq:Psi_recovery}\\[4pt]
  \hat{\boldsymbol{\Phi}}
    &= \mathbf{V}\hat{\boldsymbol{\Psi}}
       + \mathbf{1}_K\bar{\mathbf{w}}^{\!\top}
    &&\in\mathbb{R}^{K\times N},
    \label{eq:Phi_recovery}\\[4pt]
  \hat{\mathbf{B}}_z
    &= \bigl(\mathbf{C}^{\!\top}\mathbf{C}\bigr)^{-1}
       \mathbf{C}^{\!\top}\mathbf{Z}^*
    &&\in\mathbb{R}^{P\times(K{-}1)}.
    \label{eq:Bz_recovery}
\end{align}
\end{corollary}


% ---- Topic proportions -------------------------------------------
\subsection{Topic Proportions on the Simplex}
\label{subsec:proportions}

The document--topic proportions are obtained by applying the ILR
inverse~\eqref{eq:ilr_inverse} to each row of~$\mathbf{Z}^*$:
\begin{equation}
  \hat{\pi}_{ik}
  \;=\;
  \frac{\exp\!\bigl((\mathbf{V}\mathbf{z}_i^*)_k\bigr)}
       {\displaystyle\sum_{k'=1}^{K}
        \exp\!\bigl((\mathbf{V}\mathbf{z}_i^*)_{k'}\bigr)},
  \qquad i=1,\dots,M,\quad k=1,\dots,K.
  \label{eq:proportions}
\end{equation}

\begin{remark}[The closure is not a post-hoc link]
\label{rem:closure_not_posthoc}
The
mapping~\eqref{eq:proportions} is the \emph{inverse of the coordinate
system} in which the model was formulated, not an ad-hoc transformation
applied after estimation. The ILR transformation maps the simplex to
$\mathbb{R}^{K-1}$; the spectral problem is solved in
$\mathbb{R}^{K-1}$; and the closure maps back to the simplex.
This is analogous to solving a linear system in log-space and
exponentiating: the exponentiation is not an approximation but the
natural inverse of the log.
\end{remark}

% ---- Post-estimation rotation ------------------------------------
\subsection{Post-Estimation Rotation for Topic Interpretability}
\label{subsec:rotation}
 
The spectral solution of Theorem~\ref{thm:spectral} identifies the
optimal $(K{-}1)$-dimensional subspace
$\mathrm{col}(\mathbf{Z}^*)$, but the choice of basis within this
subspace is arbitrary
(Theorem~\ref{thm:identifiability_revised}\emph{(ii)}). We now show
that this rotational freedom can be exploited to improve the
interpretability of the topic--term loadings without affecting the
optimality of the solution.
 
\begin{proposition}[Rotation invariance of the profiled objective]
\label{prop:rotation_invariance}
Let $\mathbf{R}\in\mathbb{R}^{(K-1)\times(K-1)}$ be any orthogonal
matrix ($\mathbf{R}^{\!\top}\mathbf{R} = \mathbf{I}_{K-1}$). Define
the rotated parameters
\begin{equation}
  \mathbf{Z}_R = \mathbf{Z}^*\mathbf{R}, \qquad
  \boldsymbol{\Psi}_R = \mathbf{R}^{\!\top}\hat{\boldsymbol{\Psi}}, \qquad
  \mathbf{B}_{z,R} = \hat{\mathbf{B}}_z\mathbf{R}.
  \label{eq:rotated_params}
\end{equation}
Then:
\begin{enumerate}
  \item[\emph{(i)}]
    $\mathbf{Z}_R$ satisfies the constraint
    $\mathbf{Z}_R^{\!\top}\mathbf{Z}_R = \mathbf{I}_{K-1}$ and
    $\mathbf{Z}_R^{\!\top}\mathbf{1}_M = \mathbf{0}$.
  \item[\emph{(ii)}]
    The profiled objective~\eqref{eq:profiled_max} is invariant:
    $\tr(\mathbf{Z}_R^{\!\top}\mathbf{S}_z\,\mathbf{Z}_R)
     = \tr(\mathbf{Z}^{*\top}\mathbf{S}_z\,\mathbf{Z}^*)$.
  \item[\emph{(iii)}]
    The rotated parameters $(\mathbf{Z}_R,\boldsymbol{\Psi}_R,
    \mathbf{B}_{z,R},\hat{\boldsymbol{\mu}})$ achieve the same
    global optimum of~\eqref{eq:objective} as the unrotated solution.
\end{enumerate}
\end{proposition}
 
\begin{proof}
Part~\emph{(i)} follows from
$\mathbf{Z}_R^{\!\top}\mathbf{Z}_R
 = \mathbf{R}^{\!\top}\mathbf{Z}^{*\top}\mathbf{Z}^*\mathbf{R}
 = \mathbf{R}^{\!\top}\mathbf{I}_{K-1}\mathbf{R}
 = \mathbf{I}_{K-1}$
and $\mathbf{Z}_R^{\!\top}\mathbf{1}_M
 = \mathbf{R}^{\!\top}\mathbf{Z}^{*\top}\mathbf{1}_M
 = \mathbf{0}$.
Part~\emph{(ii)}: by the cyclic property of the trace,
$\tr(\mathbf{Z}_R^{\!\top}\mathbf{S}_z\,\mathbf{Z}_R)
 = \tr(\mathbf{R}^{\!\top}\mathbf{Z}^{*\top}\mathbf{S}_z\,
       \mathbf{Z}^*\mathbf{R})
 = \tr(\mathbf{Z}^{*\top}\mathbf{S}_z\,\mathbf{Z}^*
       \mathbf{R}\mathbf{R}^{\!\top})
 = \tr(\mathbf{Z}^{*\top}\mathbf{S}_z\,\mathbf{Z}^*)$.
Part~\emph{(iii)}: it suffices to verify that the rotated parameters
remain mutually consistent with the profiling
equations~\eqref{eq:Psi_hat} and~\eqref{eq:Bz_hat}. Indeed,
$K\mathbf{Z}_R^{\!\top}\widetilde{\mathbf{W}}
 = K\mathbf{R}^{\!\top}\mathbf{Z}^{*\top}\widetilde{\mathbf{W}}
 = \mathbf{R}^{\!\top}\hat{\boldsymbol{\Psi}}
 = \boldsymbol{\Psi}_R$,
and
$(\mathbf{C}^{\!\top}\mathbf{C})^{-1}\mathbf{C}^{\!\top}\mathbf{Z}_R
 = \hat{\mathbf{B}}_z\mathbf{R} = \mathbf{B}_{z,R}$.
\end{proof}
 
\noindent
Proposition~\ref{prop:rotation_invariance} establishes that the
rotation is \emph{cost-free}: the family of solutions
$\{(\mathbf{Z}^*\mathbf{R},\,\mathbf{R}^{\!\top}\hat{\boldsymbol{\Psi}},\,
\hat{\mathbf{B}}_z\mathbf{R}):\mathbf{R}\in\mathrm{O}(K{-}1)\}$ traces
out an orbit of equivalent global optima.
 
\paragraph{Effect on topic--term loadings.}
The rotation alters the topic--term matrix through the ILR deviations:
\begin{equation}
  \boldsymbol{\Phi}_R
  \;=\;
  \mathbf{V}\boldsymbol{\Psi}_R
  + \mathbf{1}_K\bar{\mathbf{w}}^{\!\top}
  \;=\;
  \mathbf{V}\mathbf{R}^{\!\top}\hat{\boldsymbol{\Psi}}
  + \mathbf{1}_K\bar{\mathbf{w}}^{\!\top}.
  \label{eq:Phi_rotated}
\end{equation}
The centroid $\bar{\mathbf{w}}$ is unchanged; only the
topic-differentiating component $\mathbf{V}\boldsymbol{\Psi}$ is
redistributed across the $K$~topics. By choosing $\mathbf{R}$
appropriately, the $K$~rows of $\boldsymbol{\Phi}_R$ can be made to
concentrate on disjoint subsets of the vocabulary, increasing topic
exclusivity and diversity.
 
\paragraph{The structural source of low exclusivity.}
To understand why rotation is needed, observe that
$\boldsymbol{\Phi} = \mathbf{V}\boldsymbol{\Psi}
+ \mathbf{1}_K\bar{\mathbf{w}}^{\!\top}$ constrains
the $K$ topic--term distributions to lie in a $(K{-}1)$-dimensional
affine subspace of $\mathbb{R}^N$ anchored at the corpus mean
$\bar{\mathbf{w}}$. The topic-differentiating part
$\mathbf{V}\boldsymbol{\Psi}$ has rank at most $K{-}1$, and its
$K$~rows are not free vectors---they are
$K$~projections of $K{-}1$ ILR directions through the contrast matrix
$\mathbf{V}$. Before rotation, these projections are aligned with the
principal eigenvectors, which maximise variance but need not produce
interpretable topic separation. Rotation within the eigenspace
reorients the ILR directions to improve separation
without leaving the optimal subspace.
 
\paragraph{Varimax rotation of ILR loadings.}
We adopt the varimax criterion~\citep{kaiser1958varimax} applied to
the transposed ILR loading matrix
$\hat{\boldsymbol{\Psi}}^{\!\top}\in\mathbb{R}^{N\times(K-1)}$,
treating it as a factor loading matrix in the standard factor-analytic
sense (terms as variables, ILR components as factors). The varimax
rotation seeks
\begin{equation}
  \mathbf{R}^*
  \;=\;
  \operatorname*{arg\,max}_{\mathbf{R}\in\mathrm{O}(K-1)}\;
  \sum_{j=1}^{K-1}
  \Biggl[
    \frac{1}{N}\sum_{n=1}^{N}
    \Bigl(\bigl(\hat{\boldsymbol{\Psi}}^{\!\top}\mathbf{R}\bigr)_{nj}
    \Bigr)^{\!4}
    \;-\;
    \Biggl(
      \frac{1}{N}\sum_{n=1}^{N}
      \Bigl(\bigl(\hat{\boldsymbol{\Psi}}^{\!\top}\mathbf{R}\bigr)_{nj}
      \Bigr)^{\!2}
    \Biggr)^{\!2}
  \Biggr],
  \label{eq:varimax}
\end{equation}
which maximises the sum of column-wise variances of the squared
loadings. This pushes each ILR component to load heavily on a small
number of terms, inducing \emph{simple structure} in the factor
pattern. Since the topic--term loadings are
$\boldsymbol{\Phi}_R = \mathbf{V}\mathbf{R}^{*\top}
\hat{\boldsymbol{\Psi}} + \mathbf{1}_K\bar{\mathbf{w}}^{\!\top}$,
simple structure in $\boldsymbol{\Psi}_R$ translates to improved
exclusivity in~$\boldsymbol{\Phi}_R$: each topic concentrates its
deviations from the corpus mean on a distinct subset of the
vocabulary.
 
\begin{remark}[Preservation of theoretical properties]
\label{rem:rotation_theory}
The rotation $\mathbf{R}^*$ is a deterministic function of
$\hat{\boldsymbol{\Psi}}$, which is itself a deterministic function
of $\mathbf{Z}^*$ and $\widetilde{\mathbf{W}}$. Since
$\mathbf{R}^*$ converges in probability to a population rotation
$\mathbf{R}_0$ (assuming the population varimax criterion has a unique
maximiser), the asymptotic results of
Section~\ref{sec:theory}---consistency
(Theorem~\ref{thm:consistency}) and asymptotic normality
(Theorem~\ref{thm:normality})---hold for the rotated path
coefficients $\mathbf{B}_{z,R} = \hat{\mathbf{B}}_z\mathbf{R}^*$
with the covariance matrix transformed as
$\boldsymbol{\Sigma}_{B_{z,R}}
 = (\mathbf{R}_0^{\!\top}\otimes\mathbf{I}_P)\,
   \boldsymbol{\Sigma}_{B_z}\,
   (\mathbf{R}_0\otimes\mathbf{I}_P)$.
The linearisation error bound
(Proposition~\ref{prop:linearisation_error}) is also
unchanged, since it depends only on $\|\mathbf{z}_i^*\|_2$, which is
invariant under orthogonal transformation:
$\|\mathbf{R}\mathbf{z}_i^*\|_2 = \|\mathbf{z}_i^*\|_2$.
\end{remark}
 
\begin{remark}[Alternative rotation criteria]
\label{rem:alternative_rotations}
While varimax on $\hat{\boldsymbol{\Psi}}^{\!\top}$ is the natural
default, other rotation criteria may be appropriate depending on the
application. Oblique rotations (e.g., promax) are not admissible in
this framework, since they would violate the orthogonality
$\mathbf{Z}_R^{\!\top}\mathbf{Z}_R = \mathbf{I}_{K-1}$; however,
any orthogonal rotation criterion (quartimax, equamax, or a custom
criterion maximising topic exclusivity directly on
$\boldsymbol{\Phi}_R$) preserves all theoretical guarantees.
\end{remark}


% ---- Efficient computation ---------------------------------------
\subsection{Efficient Computation}
\label{subsec:computation}

The matrix $\mathbf{S}_z\in\mathbb{R}^{M\times M}$ need not be formed
explicitly. Let
$\widetilde{\mathbf{W}} = \mathbf{U}\boldsymbol{\Sigma}\mathbf{V}_W^{\!\top}$
be the rank-$r$ truncated SVD and
$\mathbf{C} = \mathbf{Q}_C\mathbf{R}_C$ the thin QR decomposition.
Then $\mathbf{S}_z = \mathbf{H}\mathbf{H}^{\!\top}$, where
\begin{equation}
  \mathbf{H}
  = \bigl[\;
      \mathbf{U}\boldsymbol{\Sigma}
      \;\;\;
      \sqrt{\lambda}\,\mathbf{Q}_C
    \;\bigr]
  \;\in\;\mathbb{R}^{M\times(r+P)}.
  \label{eq:H_matrix}
\end{equation}
The top-$(K{-}1)$ left singular vectors of $\mathbf{H}$ coincide with
the desired eigenvectors of $\mathbf{S}_z$, and the computation reduces
to an eigendecomposition of the $(r{+}P)\times(r{+}P)$ Gram matrix
$\mathbf{H}^{\!\top}\mathbf{H}$, which is typically orders of magnitude
smaller than~$M\times M$.



\section{Theoretical Properties}
\label{sec:theory}

This section establishes the theoretical foundations of the ILR-Spectral-GSCA
estimator. All results concern the \emph{profiled linearised objective}
defined by~\eqref{eq:objective}--\eqref{eq:constraint}; the connection
to the original simplex problem~\eqref{eq:objective_simplex} is made
precise through the linearisation error bound in
Section~\ref{subsec:linearisation_error}.

\subsection{Standing Assumptions}
\label{subsec:assumptions}

The following assumptions are maintained throughout.

\begin{assumption}[Fixed dimensionality]
\label{ass:fixed_dim}
The number of topics $K$ and the number of covariates $P$ are fixed
as the corpus size $M\to\infty$.
\end{assumption}

\begin{assumption}[Covariate regularity]
\label{ass:covariates}
The column-centred covariate matrix
$\mathbf{C}\in\mathbb{R}^{M\times P}$ has rows
$\mathbf{c}_1,\dots,\mathbf{c}_M$ that are i.i.d.\ with
$\mathbb{E}[\mathbf{c}_i]=\mathbf{0}$,
$\mathbb{E}[\mathbf{c}_i\mathbf{c}_i^{\!\top}]
 = \boldsymbol{\Sigma}_C \succ 0$, and
$\mathbb{E}[\|\mathbf{c}_i\|^4]<\infty$.
\end{assumption}

\begin{assumption}[Spectral convergence]
\label{ass:spectral_convergence}
There exists a deterministic symmetric positive semi-definite matrix
$\boldsymbol{\Sigma}_z$, representing the population limit of the
normalised covariate-augmented similarity matrix, such that
\begin{equation}
  \bigl\|M^{-1}\mathbf{S}_z
    - \boldsymbol{\Sigma}_z\bigr\|_{\mathrm{op}}
  = O_p(M^{-1/2}).
  \label{eq:spectral_convergence}
\end{equation}
This is a high-level spectral concentration condition enabling
eigenspace perturbation arguments. It is not derived from a fully
specified generative model for~$\mathbf{W}$, but is imposed as a
standard regularity assumption for sample-to-population convergence of
the leading eigenstructure; see, e.g.,
\citet{vershynin2012close} for sufficient conditions under
sub-Gaussian designs.
\end{assumption}

\begin{assumption}[Eigengap]
\label{ass:eigengap}
The population matrix $\boldsymbol{\Sigma}_z$ has eigenvalues
$\sigma_1\geq\cdots\geq\sigma_M$ satisfying
$\delta \stackrel{\text{def}}{=}
 \sigma_{K-1} - \sigma_K > 0$;
that is, there is a strictly positive gap between the $(K{-}1)$-th
and $K$-th eigenvalues.
\end{assumption}

\begin{assumption}[Eigenvector delocalisation]
\label{ass:delocalisation}
The leading $(K{-}1)$ eigenvectors $\mathbf{u}_1,\dots,\mathbf{u}_{K-1}$
of $\mathbf{S}_z$ satisfy the incoherence condition
\begin{equation}
  \max_{1\le i\le M}\;
  \sum_{k=1}^{K-1} u_{ik}^2
  \;\leq\;
  \frac{C_0}{M}
  \label{eq:delocalisation}
\end{equation}
for some constant $C_0>0$ independent of~$M$.
Equivalently, $\max_i \|\mathbf{z}_i^*\|^2 \leq C_0/M$.
\end{assumption}

\begin{assumption}[Regularisation limit]
\label{ass:lambda}
The regularisation parameter satisfies
$\lambda = \lambda_M \to \lambda_0 \in [0,\infty)$ as $M\to\infty$.
\end{assumption}

\begin{remark}
Assumption~\ref{ass:delocalisation} is an incoherence-type regularity
condition ensuring that the leading eigenspace does not concentrate
excessively on a small number of documents. Such conditions are
standard in spectral analyses of random
matrices~\citep{candes2009exact}, although their verification depends
on the specific probabilistic structure of the data-generating
mechanism. The assumption excludes degenerate corpora in which a
small number of documents dominate the principal eigenspace.
\end{remark}

\subsection{Strengthened Linearisation Bound}
\label{subsec:strengthened_lemma}

We first establish a rigorous remainder bound for the linearisation of
the closure map (Lemma~\ref{lem:linearisation}), with an explicit
Hessian computation.

\begin{lemma}[Explicit Hessian of the closure map]
\label{lem:hessian}
Let $f(\mathbf{z}) = \closure(\exp(\mathbf{V}\mathbf{z}))$ and write
$\boldsymbol{\pi}(\mathbf{z}) = f(\mathbf{z})$.
The second derivative of the $k$-th component is
\begin{equation}
  \frac{\partial^2 f_k}{\partial z_j\,\partial z_l}
  \;=\;
  \pi_k
  \Bigl[
    (V_{kj}-\mu_j)(V_{kl}-\mu_l)
    \;-\;
    \bigl(\boldsymbol{\Sigma}_{\!\pi}^{V}\bigr)_{\!jl}
  \Bigr],
  \label{eq:hessian_fk}
\end{equation}
where $\mu_j = \sum_m\pi_m V_{mj}$ and
$(\boldsymbol{\Sigma}_{\!\pi}^{V})_{jl}
 = \sum_m\pi_m V_{mj}V_{ml} - \mu_j\mu_l$
is the $\boldsymbol{\pi}$-weighted covariance of the rows
of~$\mathbf{V}$.
Moreover, the operator norm of $D^2\!f_k(\mathbf{z})$ as a bilinear
form on $\mathbb{R}^{K-1}$ satisfies
\begin{equation}
  \|D^2\!f_k(\mathbf{z})\|_{\mathrm{op}}
  \;\leq\;
  5\,\pi_k(\mathbf{z}).
  \label{eq:hessian_bound}
\end{equation}
\end{lemma}

\begin{proof}
Starting from the first derivative
$\partial f_k/\partial z_j = \pi_k(V_{kj}-\mu_j)$,
we differentiate again:
\begin{align*}
  \frac{\partial^2 f_k}{\partial z_j\,\partial z_l}
  &= \frac{\partial\pi_k}{\partial z_l}(V_{kj}-\mu_j)
     + \pi_k\Bigl(-\frac{\partial\mu_j}{\partial z_l}\Bigr).
\end{align*}
Since $\partial\pi_k/\partial z_l = \pi_k(V_{kl}-\mu_l)$ and
\[
  \frac{\partial\mu_j}{\partial z_l}
  = \sum_m \pi_m(V_{ml}-\mu_l)V_{mj}
  = \bigl(\boldsymbol{\Sigma}_{\!\pi}^V\bigr)_{\!jl},
\]
we obtain~\eqref{eq:hessian_fk}. For the operator norm bound, let
$\mathbf{a},\mathbf{b}\in\mathbb{R}^{K-1}$ with $\|\mathbf{a}\|_2
=\|\mathbf{b}\|_2=1$. Write
$\tilde{V}_{k}^{\!\top}\mathbf{a}
= \mathbf{V}_{k\cdot}^{\!\top}\mathbf{a}
  - \boldsymbol{\mu}^{\!\top}\mathbf{a}$.
Then
\[
  \bigl|\mathbf{a}^{\!\top}D^2\!f_k\,\mathbf{b}\bigr|
  \leq \pi_k\Bigl[
    \bigl|\tilde{V}_{k}^{\!\top}\mathbf{a}\bigr|\,
    \bigl|\tilde{V}_{k}^{\!\top}\mathbf{b}\bigr|
    + \bigl\|\boldsymbol{\Sigma}_{\!\pi}^V\bigr\|_{\mathrm{op}}
  \Bigr].
\]
Since $\mathbf{V}^{\!\top}\mathbf{V}=\mathbf{I}_{K-1}$, each row of
$\mathbf{V}$ satisfies $\|\mathbf{V}_{k\cdot}\|_2\leq 1$, so
$|\mathbf{V}_{k\cdot}^{\!\top}\mathbf{a}|\leq 1$ and
$|\boldsymbol{\mu}^{\!\top}\mathbf{a}|
= |\sum_m \pi_m \mathbf{V}_{m\cdot}^{\!\top}\mathbf{a}| \leq 1$.
Therefore $|\tilde{V}_k^{\!\top}\mathbf{a}|\leq 2$ and
$|\tilde{V}_k^{\!\top}\mathbf{a}|\,|\tilde{V}_k^{\!\top}\mathbf{b}|
\leq 4$.
The matrix $\boldsymbol{\Sigma}_{\!\pi}^V$ is the
$\boldsymbol{\pi}$-weighted covariance of vectors
$\mathbf{V}_{m\cdot}$ with $\|\mathbf{V}_{m\cdot}\|_2\leq 1$,
hence $\|\boldsymbol{\Sigma}_{\!\pi}^V\|_{\mathrm{op}}
\leq \sum_m\pi_m\|\mathbf{V}_{m\cdot}\|_2^2 \leq 1$.
Combining: $|\mathbf{a}^{\!\top}D^2\!f_k\,\mathbf{b}| \leq
\pi_k(4+1) = 5\pi_k$,
which gives~\eqref{eq:hessian_bound}.
\end{proof}

\begin{proposition}[Remainder bound for the closure linearisation]
\label{prop:remainder}
For all $\mathbf{z}\in\mathbb{R}^{K-1}$ and $K\geq 2$:
\begin{equation}
  \bigl\|
    f(\mathbf{z}) - K^{-1}(\mathbf{1}_K + \mathbf{V}\mathbf{z})
  \bigr\|_2
  \;\leq\;
  \frac{5\,\exp(2\|\mathbf{z}\|_2)}{2\sqrt{K}}
  \;\|\mathbf{z}\|_2^2.
  \label{eq:remainder_bound}
\end{equation}
For $\|\mathbf{z}\|_2\leq 1$, this simplifies to
$\|R(\mathbf{z})\|_2 \leq (5e^2/2) K^{-1/2}\|\mathbf{z}\|_2^2
 < 19 K^{-1/2}\|\mathbf{z}\|_2^2$.
\end{proposition}

\begin{proof}
By the integral form of Taylor's theorem:
\[
  R_k(\mathbf{z})
  = f_k(\mathbf{z}) - f_k(\mathbf{0})
    - \mathbf{z}^{\!\top}\nabla f_k(\mathbf{0})
  = \int_0^1 (1-t)\,
    \mathbf{z}^{\!\top}D^2\!f_k(t\mathbf{z})\,\mathbf{z}\;dt.
\]
By Lemma~\ref{lem:hessian},
$|R_k(\mathbf{z})| \leq \tfrac{5}{2}\|\mathbf{z}\|_2^2\,
 \sup_{0\leq t\leq 1}\pi_k(t\mathbf{V}\mathbf{z})$.
Since each row of $\mathbf{V}$ has norm at most~1,
$|(\mathbf{V}\mathbf{z})_k|\leq\|\mathbf{z}\|_2$ for all~$k$, and
therefore
\[
  \pi_k(t\mathbf{V}\mathbf{z})
  = \frac{\exp\bigl(t(\mathbf{V}\mathbf{z})_k\bigr)}
         {\sum_{k'}\exp\bigl(t(\mathbf{V}\mathbf{z})_{k'}\bigr)}
  \leq \frac{\exp(t\|\mathbf{z}\|_2)}
            {K\exp(-t\|\mathbf{z}\|_2)}
  = \frac{\exp(2t\|\mathbf{z}\|_2)}{K}
  \leq \frac{\exp(2\|\mathbf{z}\|_2)}{K}.
\]
Thus $|R_k(\mathbf{z})|
 \leq \tfrac{5}{2}K^{-1}\exp(2\|\mathbf{z}\|_2)\|\mathbf{z}\|_2^2$
for each~$k$, and
\[
  \|R(\mathbf{z})\|_2^2
  = \sum_{k=1}^K R_k(\mathbf{z})^2
  \leq K \cdot \tfrac{25}{4}K^{-2}
       \exp(4\|\mathbf{z}\|_2)\|\mathbf{z}\|_2^4
  = \tfrac{25}{4}K^{-1}
    \exp(4\|\mathbf{z}\|_2)\|\mathbf{z}\|_2^4,
\]
giving $\|R(\mathbf{z})\|_2
 \leq \tfrac{5}{2}K^{-1/2}
      \exp(2\|\mathbf{z}\|_2)\|\mathbf{z}\|_2^2$.
\end{proof}


\subsection{Global Optimality of the Profiled Linearised Objective}
\label{subsec:global_optimality}

\begin{theorem}[Global optimum of the profiled linearised objective]
\label{thm:global}
The solution
$(\mathbf{Z}^*,\hat{\boldsymbol{\Psi}},
  \hat{\mathbf{B}}_z,\hat{\boldsymbol{\mu}})$
defined by Theorem~\ref{thm:spectral} and
Corollary~\ref{cor:recovery} is a \emph{global} minimiser
of the linearised ILR-Spectral-GSCA
objective~\eqref{eq:objective} subject
to~\eqref{eq:constraint}: no other feasible parameter set achieves a
strictly smaller value of~\eqref{eq:objective}.
\end{theorem}

\begin{proof}
For any $\mathbf{Z}$ satisfying~\eqref{eq:constraint}, the
closed-form expressions in
Propositions~\ref{prop:mu_profile}--\ref{prop:Bz_profile} are the
\emph{unique} global minimisers of~\eqref{eq:objective} over
$(\boldsymbol{\mu},\boldsymbol{\Psi},\mathbf{B}_z)$ for
fixed~$\mathbf{Z}$ (each being the solution of an unconstrained or
strictly convex quadratic). Therefore
\[
  \min_{\boldsymbol{\mu},\boldsymbol{\Psi},\mathbf{B}_z}
  \mathcal{L}(\mathbf{Z},\boldsymbol{\Psi},\mathbf{B}_z,\boldsymbol{\mu})
  = \tr(\widetilde{\mathbf{W}}^{\!\top}\widetilde{\mathbf{W}})
    - \tr(\mathbf{Z}^{\!\top}\mathbf{S}_z\,\mathbf{Z})
    + \lambda(K{-}1),
\]
and minimising over~$\mathbf{Z}$ subject
to~$\mathbf{Z}^{\!\top}\mathbf{Z}=\mathbf{I}_{K-1}$ is equivalent to
maximising $\tr(\mathbf{Z}^{\!\top}\mathbf{S}_z\,\mathbf{Z})$, which
is solved globally by the Ky~Fan theorem
(Theorem~\ref{thm:spectral}). Since each step yields a global
(not merely local) optimum, their composition is globally optimal
for~\eqref{eq:objective}.
\end{proof}

\begin{corollary}[Approximation quality for the simplex objective]
\label{cor:approx_quality}
Under Assumptions~\ref{ass:fixed_dim} and~\ref{ass:delocalisation},
let $\boldsymbol{\Theta}_{\mathrm{lin}}$ denote the linearised
topic proportions from~\eqref{eq:Theta_approx} evaluated at
$\mathbf{Z}^*$, and let
$\boldsymbol{\Theta}_{\mathrm{exact}}$ denote the exact closure
$\pi_{ik}=\closure(\exp(\mathbf{V}\mathbf{z}_i^*))_k$.
Write $\boldsymbol{\Delta}
= \boldsymbol{\Theta}_{\mathrm{exact}}-\boldsymbol{\Theta}_{\mathrm{lin}}$.
Then the \emph{per-document} discrepancy between the original simplex
objective~\eqref{eq:objective_simplex} evaluated at
$\boldsymbol{\Theta}_{\mathrm{exact}}$ and the profiled linearised
objective evaluated at $\boldsymbol{\Theta}_{\mathrm{lin}}$ satisfies
\begin{equation}
  \frac{1}{M}
  \bigl|
    \mathcal{L}(\boldsymbol{\Theta}_{\mathrm{exact}},
      \hat{\boldsymbol{\Phi}},\hat{\mathbf{B}})
    - \mathcal{L}_{\mathrm{lin}}(\mathbf{Z}^*,
      \hat{\boldsymbol{\Psi}},\hat{\mathbf{B}}_z,\hat{\boldsymbol{\mu}})
  \bigr|
  \;\leq\;
  \frac{2}{M}\bigl(\|\widetilde{\mathbf{W}}\|_F
         \|\hat{\boldsymbol{\Phi}}\|_{\mathrm{op}}
         + \lambda\bigr)
  \|\boldsymbol{\Delta}\|_F
  + \frac{1+\lambda}{M}\|\boldsymbol{\Delta}\|_F^2
  = O_p(M^{-1}).
  \label{eq:approx_quality}
\end{equation}
That is, the average objective discrepancy per document vanishes at
rate $O_p(M^{-1})$.
\end{corollary}

\begin{proof}
The objective~\eqref{eq:objective_simplex} is a sum of two squared
Frobenius norms. For any matrices $\mathbf{A}$, $\mathbf{B}$ and
perturbation $\boldsymbol{\Delta}$:
$\bigl|\|\mathbf{A}-\mathbf{B}-\boldsymbol{\Delta}\|_F^2
 - \|\mathbf{A}-\mathbf{B}\|_F^2\bigr|
\leq 2\|\mathbf{A}-\mathbf{B}\|_F\|\boldsymbol{\Delta}\|_F
    + \|\boldsymbol{\Delta}\|_F^2$.
Applying this to both terms of~\eqref{eq:objective_simplex}:
the reconstruction term contributes at most
$2\|\widetilde{\mathbf{W}}\|_F
 \|\hat{\boldsymbol{\Phi}}\|_{\mathrm{op}}
 \|\boldsymbol{\Delta}\|_F
 + \|\hat{\boldsymbol{\Phi}}\|_{\mathrm{op}}^2
   \|\boldsymbol{\Delta}\|_F^2$,
and the structural term at most
$2\lambda\|\boldsymbol{\Delta}\|_F
 + \lambda\|\boldsymbol{\Delta}\|_F^2$.
Combining, dividing by~$M$, and using
Proposition~\ref{prop:linearisation_error}
($\|\boldsymbol{\Delta}\|_F^2 = O(M^{-1})$, so
$\|\boldsymbol{\Delta}\|_F = O(M^{-1/2})$):
the leading term scales as
$M^{-1}\cdot O_p(\sqrt{M})\cdot O(M^{-1/2}) = O_p(M^{-1})$,
where $M^{-1}\|\widetilde{\mathbf{W}}\|_F = O_p(M^{-1/2})$
follows from bounded per-document moment conditions.
The quadratic term $M^{-1}\|\boldsymbol{\Delta}\|_F^2 = O(M^{-2})$
is of smaller order, giving the stated rate.
\end{proof}


\subsection{Identifiability}
\label{subsec:identifiability}

\begin{theorem}[Identifiability of the spectral solution]
\label{thm:identifiability_revised}
Under Assumption~\ref{ass:eigengap}:
\begin{enumerate}
  \item[\emph{(i)}]
    If the top $K{-}1$ eigenvalues of $\mathbf{S}_z$ are \emph{simple}
    (i.e.,\ $\sigma_1>\sigma_2>\cdots>\sigma_{K-1}>\sigma_K$), then
    $\mathbf{Z}^*$ is determined up to sign flips of individual columns.
    The full parameter set
    $(\hat{\boldsymbol{\Psi}},\hat{\mathbf{B}}_z,\hat{\boldsymbol{\mu}})$
    is identified up to the same sign convention.

  \item[\emph{(ii)}]
    If some top eigenvalues coincide, only the subspace
    $\mathrm{col}(\mathbf{Z}^*)\subset\mathbb{R}^M$ is identified.
    The individual columns of $\mathbf{Z}^*$ are determined only up to
    a $(K{-}1)\times(K{-}1)$ orthogonal transformation within the
    eigenspace.
\end{enumerate}
\end{theorem}

\begin{proof}
Part~\emph{(i)}: eigenvectors corresponding to simple eigenvalues of a
real symmetric matrix are unique up to scalar multiplication; the
unit-norm constraint restricts this to a sign choice.
Since $\hat{\boldsymbol{\Psi}} = K\mathbf{Z}^{*\top}
\widetilde{\mathbf{W}}$ and $\hat{\mathbf{B}}_z
= (\mathbf{C}^{\!\top}\mathbf{C})^{-1}\mathbf{C}^{\!\top}\mathbf{Z}^*$
are linear in~$\mathbf{Z}^*$, the sign ambiguity propagates directly.

Part~\emph{(ii)}: when $\sigma_j=\sigma_{j+1}$ for some
$j\in\{1,\dots,K{-}2\}$, any orthonormal basis of the corresponding
eigenspace solves~\eqref{eq:profiled_max}. The identified object is
the orthogonal projector
$\mathbf{Z}^*\mathbf{Z}^{*\top}$, which determines the subspace but
not the basis.
\end{proof}

\begin{remark}[Resolving sign and rotational ambiguity]
\label{rem:resolving_ambiguity}
In case~\emph{(i)}, column signs can be fixed by a deterministic
convention, e.g., requiring that the entry of $\hat{\boldsymbol{\Psi}}_k$
with largest absolute value is positive. In case~\emph{(ii)}, a rotation
criterion (e.g., varimax on $\hat{\boldsymbol{\Psi}}$) selects a
specific basis within the eigenspace for interpretability. Such post-hoc
orientation is standard in factor analysis and does not affect the
statistical content of the solution.
\end{remark}


\subsection{Consistency of the Path Coefficients}
\label{subsec:consistency}

Let $\mathbf{U}_0 = [\mathbf{u}_{01},\dots,\mathbf{u}_{0,K-1}]$
denote the matrix whose columns are the leading $K{-}1$ eigenvectors
of the population matrix $\boldsymbol{\Sigma}_z$
(Assumption~\ref{ass:spectral_convergence}), aligned according to the
identification convention adopted in
Theorem~\ref{thm:identifiability_revised}. Let
$\mathbf{u}_{0i}^{\!\top}\in\mathbb{R}^{1\times(K-1)}$ denote its
$i$-th row. We define the \emph{population path coefficient matrix}
\begin{equation}
  \mathbf{B}_{z,0}
  \;\stackrel{\text{def}}{=}\;
  \boldsymbol{\Sigma}_C^{-1}\,
  \mathbb{E}\!\bigl[\mathbf{c}_i\,\mathbf{u}_{0i}^{\!\top}\bigr]
  \;\in\;\mathbb{R}^{P\times(K-1)},
  \label{eq:Bz0_def}
\end{equation}
which is the regression of the population ILR scores on the covariates.

\begin{theorem}[Consistency of~$\hat{\mathbf{B}}_z$]
\label{thm:consistency}
Under
Assumptions~\ref{ass:fixed_dim}--\ref{ass:lambda}:
\[
  \hat{\mathbf{B}}_z
  \;\xrightarrow{p}\;
  \mathbf{B}_{z,0}.
\]
\end{theorem}

\begin{proof}
The proof proceeds in three steps.

\emph{Step~1 (covariate Gram matrix).}
By the law of large numbers and
Assumption~\ref{ass:covariates}:
\begin{equation}
  M^{-1}\mathbf{C}^{\!\top}\mathbf{C}
  \;\xrightarrow{p}\;
  \boldsymbol{\Sigma}_C \succ 0.
  \label{eq:gram_lln}
\end{equation}

\emph{Step~2 (eigenvector convergence via Davis--Kahan).}
Let $\hat{\mathbf{U}} = [\mathbf{u}_1,\dots,\mathbf{u}_{K-1}]$ and
$\mathbf{U}_0 = [\mathbf{u}_{01},\dots,\mathbf{u}_{0,K-1}]$ denote the
sample and population leading eigenvector matrices.
By the Davis--Kahan $\sin\Theta$
theorem~\citep{davis1970rotation}:
\begin{equation}
  \bigl\|
    \sin\Theta(\hat{\mathbf{U}},\mathbf{U}_0)
  \bigr\|_{\mathrm{op}}
  \;\leq\;
  \frac{\|M^{-1}\mathbf{S}_z - \boldsymbol{\Sigma}_z\|_{\mathrm{op}}}
       {\delta}
  \;=\;
  O_p(M^{-1/2}),
  \label{eq:davis_kahan}
\end{equation}
where $\delta>0$ is the eigengap from
Assumption~\ref{ass:eigengap} and the rate follows from
Assumption~\ref{ass:spectral_convergence}.
Consequently, there exists a sequence of $(K{-}1)\times(K{-}1)$
orthogonal matrices $\mathbf{W}_M$ (resolving the sign/rotation
ambiguity) such that
\begin{equation}
  \bigl\|\hat{\mathbf{U}} - \mathbf{U}_0\mathbf{W}_M\bigr\|_F
  = O_p(M^{-1/2}).
  \label{eq:eigvec_convergence}
\end{equation}

\emph{Step~3 (continuous mapping).}
The estimator (absorbing the orthogonal alignment $\mathbf{W}_M$) is
\[
  \hat{\mathbf{B}}_z
  = \bigl(M^{-1}\mathbf{C}^{\!\top}\mathbf{C}\bigr)^{-1}
    \bigl(M^{-1}\mathbf{C}^{\!\top}\hat{\mathbf{U}}\bigr).
\]
By~\eqref{eq:gram_lln}, the first factor converges to
$\boldsymbol{\Sigma}_C^{-1}$.
For the second factor:
\[
  M^{-1}\mathbf{C}^{\!\top}\hat{\mathbf{U}}
  = M^{-1}\mathbf{C}^{\!\top}\mathbf{U}_0\mathbf{W}_M
    + M^{-1}\mathbf{C}^{\!\top}
      (\hat{\mathbf{U}}-\mathbf{U}_0\mathbf{W}_M).
\]
The first term converges to
$\mathbb{E}[\mathbf{c}_i\mathbf{u}_{0i}^{\!\top}]\mathbf{W}_M$ by
the LLN. The second term is bounded by
$\|M^{-1}\mathbf{C}^{\!\top}\|_{\mathrm{op}}
 \|\hat{\mathbf{U}}-\mathbf{U}_0\mathbf{W}_M\|_F
 = O_p(1)\cdot O_p(M^{-1/2}) = o_p(1)$.
The continuous mapping theorem for the map
$(\mathbf{A},\mathbf{D})\mapsto\mathbf{A}^{-1}\mathbf{D}$ (continuous
at $(\boldsymbol{\Sigma}_C, \mathbb{E}[\mathbf{c}_i\mathbf{u}_{0i}^{\!\top}])$
since $\boldsymbol{\Sigma}_C\succ 0$) yields the stated convergence.
\end{proof}


\subsection{Asymptotic Normality of the Path Coefficients}
\label{subsec:asymptotic_normality}

To establish the asymptotic distribution of $\hat{\mathbf{B}}_z$, we
require a first-order expansion of the sample eigenvectors.

\begin{lemma}[Eigenvector perturbation expansion]
\label{lem:eigvec_expansion}
Under Assumptions~\ref{ass:spectral_convergence}
and~\ref{ass:eigengap}, for each $k=1,\dots,K{-}1$:
\begin{equation}
  \mathbf{u}_k - \mathbf{u}_{0k}
  \;=\;
  \sum_{\ell\neq k}
  \frac{\mathbf{u}_{0\ell}^{\!\top}
        (M^{-1}\mathbf{S}_z - \boldsymbol{\Sigma}_z)\,
        \mathbf{u}_{0k}}
       {\sigma_k - \sigma_\ell}\;
  \mathbf{u}_{0\ell}
  \;+\;
  \mathbf{r}_{k,M},
  \label{eq:eigvec_expansion}
\end{equation}
where $\|\mathbf{r}_{k,M}\| = O_p(M^{-1})$.
\end{lemma}

\begin{proof}
This is the classical first-order perturbation expansion for
eigenvectors of symmetric matrices under simple eigenvalues. For the
general theory, see \citet{kato1995perturbation}~(Chapter~II,
Theorem~2.6 and the discussion in \S6.2);
for the statistical formulation with remainder control, see
\citet{anderson1963asymptotic}~(Theorem~1)
and~\citet{tyler1981asymptotic}~(Theorem~2.1).
Let $\mathbf{E} = M^{-1}\mathbf{S}_z - \boldsymbol{\Sigma}_z$ with
$\|\mathbf{E}\|_{\mathrm{op}} = O_p(M^{-1/2})$ by
Assumption~\ref{ass:spectral_convergence}.
The resolvent identity applied to the perturbed eigenvalue problem
$(\boldsymbol{\Sigma}_z + \mathbf{E})\mathbf{u}_k
= \hat\sigma_k\mathbf{u}_k$ yields
\[
  \mathbf{u}_k
  = \mathbf{u}_{0k}
    + (\sigma_k\mathbf{I} - \boldsymbol{\Sigma}_z)^{\dagger}
      \mathbf{E}\,\mathbf{u}_{0k}
    + O_p(\|\mathbf{E}\|^2),
\]
where $(\cdot)^{\dagger}$ denotes the pseudoinverse restricted to the
orthogonal complement of $\mathbf{u}_{0k}$.
Expanding the pseudoinverse in the eigenbasis of $\boldsymbol{\Sigma}_z$
gives~\eqref{eq:eigvec_expansion} with remainder
$\|\mathbf{r}_{k,M}\| = O(\|\mathbf{E}\|^2_{\mathrm{op}})
= O_p(M^{-1})$.
\end{proof}

\begin{theorem}[First-order asymptotic normality of~$\hat{\mathbf{B}}_z$]
\label{thm:normality}
Under Assumptions~\ref{ass:fixed_dim}--\ref{ass:lambda}, assuming
additionally that the top $K{-}1$ eigenvalues of
$\boldsymbol{\Sigma}_z$ are simple, the estimator
$\hat{\mathbf{B}}_z$ admits the following first-order asymptotic
normal characterisation based on the eigenspace perturbation expansion
of Lemma~\ref{lem:eigvec_expansion}:
\begin{equation}
  \sqrt{M}\,\vect\!\bigl(\hat{\mathbf{B}}_z - \mathbf{B}_{z,0}\bigr)
  \;\xrightarrow{d}\;
  \mathcal{N}\!\bigl(\mathbf{0},\;\boldsymbol{\Sigma}_{B_z}\bigr),
  \label{eq:Bz_clt}
\end{equation}
where the first-order asymptotic covariance matrix
$\boldsymbol{\Sigma}_{B_z}$ is given by
\begin{equation}
  \boldsymbol{\Sigma}_{B_z}
  \;=\;
  \bigl(\mathbf{I}_{K-1}\otimes\boldsymbol{\Sigma}_C^{-1}\bigr)
  \;\boldsymbol{\Gamma}\;
  \bigl(\mathbf{I}_{K-1}\otimes\boldsymbol{\Sigma}_C^{-1}\bigr),
  \label{eq:Sigma_Bz}
\end{equation}
with $\boldsymbol{\Gamma}\in\mathbb{R}^{P(K-1)\times P(K-1)}$
defined blockwise: for $k,k'=1,\dots,K{-}1$,
\begin{equation}
  \boldsymbol{\Gamma}_{kk'}
  \;=\;
  \sum_{\ell:\,\ell\neq k}\;
  \sum_{\ell':\,\ell'\neq k'}
  \frac{1}{(\sigma_k-\sigma_\ell)(\sigma_{k'}-\sigma_{\ell'})}
  \;\mathbb{E}\!\bigl[
    \mathbf{c}_i\mathbf{c}_i^{\!\top}\,
    (\mathbf{u}_{0\ell}^{\!\top}\boldsymbol{\xi}_{ik})
    (\mathbf{u}_{0\ell'}^{\!\top}\boldsymbol{\xi}_{ik'})
  \bigr],
  \label{eq:Gamma_block}
\end{equation}
where $\boldsymbol{\xi}_{ik}$ is the centred contribution of
document~$i$ to the perturbation of the $k$-th eigenvector, defined as
follows. Let $\widetilde{\mathbf{w}}_i\in\mathbb{R}^N$ denote the
centred term-frequency vector for document~$i$ (the $i$-th row
of~$\widetilde{\mathbf{W}}$, viewed as a column vector). Then
\begin{equation}
  \boldsymbol{\xi}_{ik}
  \;=\;
  \bigl(
    \widetilde{\mathbf{w}}_i\widetilde{\mathbf{w}}_i^{\!\top}
    - \mathbb{E}\bigl[
        \widetilde{\mathbf{w}}\widetilde{\mathbf{w}}^{\!\top}
      \bigr]
    + \lambda_0\bigl(
        \mathbf{c}_i\mathbf{c}_i^{\!\top}
        - \boldsymbol{\Sigma}_C
      \bigr)
  \bigr)\,\mathbf{u}_{0k},
  \label{eq:xi_ik}
\end{equation}
i.e., the action of the centred rank-one perturbation of $M^{-1}\mathbf{S}_z$
on the population eigenvector $\mathbf{u}_{0k}$.
\end{theorem}

\begin{proof}
Substituting the eigenvector
expansion~\eqref{eq:eigvec_expansion} into
$\hat{\mathbf{B}}_z = (M^{-1}\mathbf{C}^{\!\top}\mathbf{C})^{-1}
(M^{-1}\mathbf{C}^{\!\top}\hat{\mathbf{U}})$, the $k$-th column of
$\hat{\mathbf{B}}_z - \mathbf{B}_{z,0}$ equals
\begin{align}
  \hat{\mathbf{b}}_{z,k} - \mathbf{b}_{z,0k}
  &= \boldsymbol{\Sigma}_C^{-1}
     \cdot M^{-1}\mathbf{C}^{\!\top}
     \Bigl[
       \sum_{\ell\neq k}
       \frac{\mathbf{u}_{0\ell}^{\!\top}\mathbf{E}\,\mathbf{u}_{0k}}
            {\sigma_k - \sigma_\ell}\,\mathbf{u}_{0\ell}
     \Bigr]
     + o_p(M^{-1/2}) \notag\\
  &= \boldsymbol{\Sigma}_C^{-1}
     \sum_{\ell\neq k}
     \frac{1}{\sigma_k - \sigma_\ell}
     \Bigl(
       M^{-1}\sum_{i=1}^M
       \mathbf{c}_i\,
       (\mathbf{u}_{0\ell}^{\!\top}\boldsymbol{\xi}_{ik})
     \Bigr)
     + o_p(M^{-1/2}).
  \label{eq:Bz_expansion}
\end{align}
The summands
$\mathbf{c}_i(\mathbf{u}_{0\ell}^{\!\top}\boldsymbol{\xi}_{ik})$
are centred (since $\mathbb{E}[\boldsymbol{\xi}_{ik}]=\mathbf{0}$
by construction) and have finite
second moments under
Assumptions~\ref{ass:covariates}--\ref{ass:spectral_convergence}.
Although the sample eigenvectors are global objects depending on all
$M$~documents, the first-order perturbation expansion rewrites their
leading fluctuation as a linear functional of the empirical perturbation
$\mathbf{E} = M^{-1}\mathbf{S}_z - \boldsymbol{\Sigma}_z
 = M^{-1}\sum_{i=1}^M \boldsymbol{\xi}_{ik}\mathbf{u}_{0k}^{\!\top}
   + \cdots$,
which is itself an average of document-level rank-one contributions.
Under the i.i.d.\ sampling of documents
(Assumption~\ref{ass:covariates}), the first-order terms
$\mathbf{c}_i(\mathbf{u}_{0\ell}^{\!\top}\boldsymbol{\xi}_{ik})$ are
therefore i.i.d.\ across~$i$, which justifies the multivariate central
limit argument.
The multivariate central limit theorem gives
\[
  \sqrt{M}\sum_{\ell\neq k}
  \frac{1}{\sigma_k-\sigma_\ell}
  \Bigl(
    M^{-1}\textstyle\sum_i
    \mathbf{c}_i\,\mathbf{u}_{0\ell}^{\!\top}\boldsymbol{\xi}_{ik}
  \Bigr)
  \xrightarrow{d}
  \mathcal{N}\bigl(\mathbf{0},\,
    \boldsymbol{\Gamma}_{kk}\bigr),
\]
and the cross-topic covariance $\boldsymbol{\Gamma}_{kk'}$ arises from
the joint CLT over all $k,k'=1,\dots,K{-}1$.
Applying the Cram\'er--Wold device and Slutsky's theorem (for the
replacement of $\boldsymbol{\Sigma}_C^{-1}$) yields~\eqref{eq:Bz_clt}
with covariance~\eqref{eq:Sigma_Bz}--\eqref{eq:Gamma_block}.
\end{proof}

\begin{remark}[First-order covariance structure]
\label{rem:no_bootstrap}
The covariance matrix $\boldsymbol{\Sigma}_{B_z}$
in~\eqref{eq:Sigma_Bz}--\eqref{eq:Gamma_block} is induced by the
joint first-order perturbation law of the leading eigenspace and
the regression step. It involves the population eigenvalues
$\sigma_k$, the population eigenvectors $\mathbf{u}_{0k}$, and
moments of the data, all of which can be consistently estimated by
plug-in substitution of their sample counterparts.
The resulting sandwich estimator
$\hat{\boldsymbol{\Sigma}}_{B_z}$ provides asymptotically valid
standard errors and confidence intervals for the entries
of~$\mathbf{B}_z$ under the stated regularity conditions.
The qualifier ``first-order'' refers to the fact that the
covariance~\eqref{eq:Gamma_block} is derived from the linear term
of the eigenvector perturbation expansion~\eqref{eq:eigvec_expansion};
the $O_p(M^{-1})$ remainder contributes only at higher order.
This analytical approach contrasts with the GSCA-TM two-stage
estimator, which requires a parametric bootstrap to propagate
first-stage uncertainty.
\end{remark}


\subsection{Topic Proportions via the Delta Method}
\label{subsec:delta_method}

\begin{remark}[Asymptotic distribution of topic proportions]
\label{rem:pi_delta_method}
Once the asymptotic linearity of the eigenvector
estimates~\eqref{eq:eigvec_expansion} is established, the delta method
can be applied to any sufficiently smooth functional of the ILR topic
scores. In particular, for \emph{aggregate} functionals of the form
$M^{-1}\sum_i h(\hat{\boldsymbol{\pi}}_i)$ (e.g., corpus-level topic
prevalences), the delta method combined with the CLT
of Theorem~\ref{thm:normality} yields asymptotic normality with a
covariance derived from the Jacobian of the ILR inverse,
$\mathbf{J}_f(\mathbf{z}) = K^{-1}[\diag(\boldsymbol{\pi})
- \boldsymbol{\pi}\boldsymbol{\pi}^{\!\top}]\mathbf{V}$, composed with
the eigenvector perturbation. Asymptotic statements for \emph{individual}
document proportions $\hat{\boldsymbol{\pi}}_i$ require careful treatment
of the dependence between the $i$-th score and the remaining documents
through the eigenvector structure, and are deferred to future work.
\end{remark}


\subsection{Linearisation Error and Connection to the Simplex Objective}
\label{subsec:linearisation_error}

\begin{proposition}[Aggregate linearisation error]
\label{prop:linearisation_error}
Under Assumptions~\ref{ass:fixed_dim} and~\ref{ass:delocalisation},
let $\boldsymbol{\Theta}_{\mathrm{exact}}$ have rows
$\closure(\exp(\mathbf{V}\mathbf{z}_i^*))^{\!\top}$ and let
$\boldsymbol{\Theta}_{\mathrm{lin}}
= K^{-1}(\mathbf{1}_M\mathbf{1}_K^{\!\top}
+ \mathbf{Z}^*\mathbf{V}^{\!\top})$.
Then
\begin{equation}
  \frac{1}{MK}\,
  \bigl\|\boldsymbol{\Theta}_{\mathrm{exact}}
    - \boldsymbol{\Theta}_{\mathrm{lin}}\bigr\|_F^2
  \;\leq\;
  \frac{25\,C_0(K{-}1)\,e^{4\sqrt{C_0/M}}}{4K^2 M^2}
  \;=\;
  O\!\bigl(M^{-2}\bigr).
  \label{eq:linearisation_error}
\end{equation}
In particular, the mean squared entry-wise discrepancy between the exact
and linearised topic proportions vanishes at rate $O(M^{-2})$.
\end{proposition}

\begin{proof}
By Proposition~\ref{prop:remainder},
$\|R(\mathbf{z}_i^*)\|_2^2
 \leq \tfrac{25}{4}K^{-1}
      \exp(4\|\mathbf{z}_i^*\|_2)\|\mathbf{z}_i^*\|_2^4$.
Using Assumption~\ref{ass:delocalisation},
$\|\mathbf{z}_i^*\|_2 \leq \sqrt{C_0/M}$ for all~$i$, and the constraint
$\sum_i\|\mathbf{z}_i^*\|_2^2 = K{-}1$ implies
$\sum_i\|\mathbf{z}_i^*\|_2^4
 \leq (\max_i\|\mathbf{z}_i^*\|_2^2)(K{-}1)
 \leq C_0(K{-}1)/M$.
Therefore
\[
  \frac{1}{MK}\sum_{i=1}^M \|R(\mathbf{z}_i^*)\|_2^2
  \leq \frac{25\,\exp(4\sqrt{C_0/M})}{4MK^2}
       \cdot\frac{C_0(K{-}1)}{M}
  = \frac{25\,C_0(K{-}1)\,e^{4\sqrt{C_0/M}}}{4K^2 M^2},
\]
which is $O(M^{-2})$ for $K$ fixed.
\end{proof}

\subsection{Scope of the Theoretical Guarantees}
\label{subsec:scope}

To aid the reader, we summarise the status of each result and
delineate the inferential perimeter:
%
\begin{enumerate}
  \item \textbf{Exact results for the profiled linearised objective.}
    Global optimality (Theorem~\ref{thm:global}),
    identifiability (Theorem~\ref{thm:identifiability_revised}), and
    the spectral characterisation (Theorem~\ref{thm:spectral})
    hold without approximation for the ILR-linearised
    problem~\eqref{eq:objective}--\eqref{eq:constraint}.

  \item \textbf{Asymptotic results under standing assumptions.}
    Consistency (Theorem~\ref{thm:consistency}) and first-order
    asymptotic normality (Theorem~\ref{thm:normality}) of
    $\hat{\mathbf{B}}_z$ hold under
    Assumptions~\ref{ass:fixed_dim}--\ref{ass:lambda} as
    $M\to\infty$. The inferential theory for $\mathbf{B}_z$ is
    closed: standard errors and confidence intervals can be
    constructed from the analytical covariance formula without
    bootstrap.

  \item \textbf{Approximation guarantee.}
    The connection to the original simplex
    objective~\eqref{eq:objective_simplex} is mediated by the
    linearisation error bound
    (Proposition~\ref{prop:linearisation_error} and
    Corollary~\ref{cor:approx_quality}), which vanishes at rate
    $O(M^{-2})$ under the delocalisation
    condition~(Assumption~\ref{ass:delocalisation}).

  \item \textbf{Open extensions.}
    Extending pointwise asymptotic inference to individual topic
    compositions $\hat{\boldsymbol{\pi}}_i$ requires additional
    control of the dependence induced by the global spectral
    normalisation and is deferred to future work.
    The delta method applies readily to \emph{aggregate} functionals
    of topic proportions (Remark~\ref{rem:pi_delta_method}).
\end{enumerate}

\noindent
In summary, the current theory is exact for the profiled linearised
objective, asymptotically closed at the first-order perturbation level
for the structural coefficient matrix~$\mathbf{B}_z$ under the stated
regularity conditions, and only partially developed for
document-specific topic proportions, whose pointwise inference is left
to future work.

\noindent
In summary, the current theory is exact for the profiled linearised
objective, asymptotically closed at the first-order perturbation level
for the structural coefficient matrix $\mathbf{B}_z$ under the stated
regularity conditions, and only partially developed for
document-specific topic proportions, whose pointwise inference is left
to future work.


% ====================================================================
\section{Connection to Existing Frameworks}
\label{sec:connections}
% ====================================================================

The ILR-Spectral-GSCA topic model sits at the intersection of three established
frameworks:
%
\begin{enumerate}
  \item \textbf{Topic modeling.}
    For $\lambda=0$, the solution reduces to the singular value
    decomposition of the centred document--term matrix---equivalent
    to Latent Semantic Analysis (LSA) applied to centred data, with
    topic proportions obtained via ILR back-transformation instead of
    arbitrary thresholding.

  \item \textbf{Structural equation modeling.}
    The structural penalty
    $\lambda\|\mathbf{Z}-\mathbf{C}\mathbf{B}_z\|_F^2$ is the standard
    GSCA measurement model~\citep{hwang2004generalized}, with
    $\mathbf{Z}$ playing the role of latent component scores.
    The contribution here is that the latent space is not arbitrary but
    is the ILR transform of the probability simplex, giving the scores
    a compositional interpretation.

  \item \textbf{Compositional data analysis.}
    The estimation procedure in ILR coordinates respects the Aitchison
    geometry of the simplex~\citep{aitchison1986statistical}. The
    structural penalty
    $\|\mathbf{Z}-\mathbf{C}\mathbf{B}_z\|_F^2$ in ILR coordinates
    equals the sum of squared Aitchison distances between the observed
    topic compositions and the covariate-predicted compositions
    (Remark~\ref{rem:what_reparam_achieves}); this compositional
    interpretation pertains to the reparameterised objective, not to
    the original Frobenius penalty in $\boldsymbol{\Theta}$-space.
    The centering $\widetilde{\mathbf{W}} = \mathbf{H}_M\mathbf{W}$
    removes the compositional centroid, and the ILR loadings
    $\hat{\boldsymbol{\Psi}}$ describe log-ratio contrasts between
    topics---the natural ``effects'' in compositional
    analysis~\citep{egozcue2003isometric,
    pawlowskyglahn2015modeling}.
\end{enumerate}


% ====================================================================
\section{Algorithm}
\label{sec:algorithm}
% ====================================================================

\begin{algorithm}[t]
\caption{ILR-Spectral-GSCA Structural Topic Model}
\label{(alg:ilr_spectral_gsca)}
\begin{algorithmic}[1]
\REQUIRE Document--term matrix
  $\mathbf{W}\in\mathbb{R}^{M\times N}_{\geq 0}$,
  column-centred covariate matrix
  $\mathbf{C}\in\mathbb{R}^{M\times P}$,
  number of topics $K$,
  regularisation $\lambda>0$,
  ILR contrast matrix
  $\mathbf{V}\in\mathbb{R}^{K\times(K-1)}$.
\ENSURE Topic proportions $\hat{\boldsymbol{\Pi}}$,
  topic--term loadings $\hat{\boldsymbol{\Phi}}$,
  ILR path coefficients $\hat{\mathbf{B}}_z$.

\medskip
\STATE \textbf{Centre.}\;
  $\widetilde{\mathbf{W}} \leftarrow \mathbf{W}
    - M^{-1}\mathbf{1}_M\mathbf{1}_M^{\!\top}\mathbf{W}$
\STATE \textbf{Thin SVD.}\;
  $\mathbf{U}\boldsymbol{\Sigma}\mathbf{V}_W^{\!\top}
    \leftarrow \text{SVD}_r(\widetilde{\mathbf{W}})$
\STATE \textbf{Covariate QR.}\;
  $\mathbf{Q}_C\mathbf{R}_C \leftarrow \text{QR}(\mathbf{C})$
\STATE \textbf{Augmented factor.}\;
  $\mathbf{H} \leftarrow [\,\mathbf{U}\boldsymbol{\Sigma}
    \;\;\;\sqrt{\lambda}\,\mathbf{Q}_C\,]$
\STATE \textbf{Reduced eigendecomposition.}\;
  $\mathbf{H}^{\!\top}\mathbf{H}
    = \mathbf{E}\,\diag(\sigma_1,\dots,\sigma_{r+P})\,
      \mathbf{E}^{\!\top}$
\STATE \textbf{ILR topic scores.}\;
  $\mathbf{Z}^* \leftarrow \mathbf{H}\mathbf{E}_{1:K-1}\,
    \diag(\sigma_1^{-1/2},\dots,\sigma_{K-1}^{-1/2})$
\STATE \textbf{ILR loadings.}\;
  $\hat{\boldsymbol{\Psi}} \leftarrow
    K\,\mathbf{Z}^{*\top}\widetilde{\mathbf{W}}$
\STATE \textbf{Varimax rotation (optional).}\;
  $\mathbf{R}^* \leftarrow \text{varimax}(\hat{\boldsymbol{\Psi}}^{\!\top})$;\;
  $\mathbf{Z}^* \leftarrow \mathbf{Z}^*\mathbf{R}^*$;\;
  $\hat{\boldsymbol{\Psi}} \leftarrow \mathbf{R}^{*\top}\hat{\boldsymbol{\Psi}}$;\;
  $\hat{\mathbf{B}}_z \leftarrow \hat{\mathbf{B}}_z\mathbf{R}^*$
\STATE \textbf{Topic--term matrix.}\;
  $\hat{\boldsymbol{\Phi}} \leftarrow
    \mathbf{V}\hat{\boldsymbol{\Psi}}
    + \mathbf{1}_K\bar{\mathbf{w}}^{\!\top}$
\STATE \textbf{Path coefficients.}\;
  $\hat{\mathbf{B}}_z \leftarrow
    (\mathbf{C}^{\!\top}\mathbf{C})^{-1}
    \mathbf{C}^{\!\top}\mathbf{Z}^*$
\STATE \textbf{Topic proportions.}\;
  $\hat{\pi}_{ik} \leftarrow
    \exp((\mathbf{V}\mathbf{z}_i^*)_k)\big/
    \sum_{k'}\exp((\mathbf{V}\mathbf{z}_i^*)_{k'})$
  \quad$\forall\; i,k$

\RETURN $\hat{\boldsymbol{\Pi}},\;
  \hat{\boldsymbol{\Phi}},\;
  \hat{\mathbf{B}}_z$
\end{algorithmic}
\end{algorithm}

\noindent
The algorithm is \emph{non-iterative}: it consists of a fixed sequence
of matrix decompositions and linear-algebra operations, with no
convergence loop and no random initialisation. Its output is
deterministic for any given input.


\bibliographystyle{apalike}
\bibliography{bibliography}
\end{document}